\begin{document}
\maketitle

\begin{abstract}
We present the complete axiomatic foundation of Topological Fixed Point Theory (TFPT), 
which derives the fine-structure constant $\alpha$, within the declared scheme and renormalization condition, and other fundamental quantities 
from eight minimal assumptions rooted in topology, geometry, and quantum field theory. 
The theory rests on two fundamental invariants: the topological coupling 
$\cthree = 1/(8\pi)$ from 11D Chern--Simons quantization and the geometric scale 
$\phiz = 1/(6\pi) + 3/(256\pi^4)$ from M\"obius fiber geometry. 
These invariants, combined with the Standard Model abelian trace $\bone = 41/10$,
yield a cubic fixed-point equation (CFE) whose unique positive root is fixed by geometric self-consistency on the orientable double cover ($k=2$).
A two-defect correction derived from SU(5) holonomy degeneracy ($g=3+2=5$) yields $\delta_2=\frac54\,(\delta_{\mathrm{top}})^2$ and a kernel prediction
$\alpha^{-1}(0)=\num[round-mode=places,round-precision=12]{\TFPTalphaInvGOverFourNum}$ (CODATA $z=\num[round-mode=places,round-precision=3]{\TFPTalphaZGOverFourNum}$).
The one-defect truncation baseline yields $\alpha^{-1}(0)=\num[round-mode=places,round-precision=10]{\TFPTalphaInvSelfNum}$ (a $-0.037$ ppm deviation) and is reported only for series control and sensitivity diagnostics.
Under the declared $\overline{\mathrm{MS}}$-at-$M_Z$ comparison policy, the primary comparison observable is
$\overline{\alpha}^{(5)}(M_Z)^{-1}=\num[round-mode=places,round-precision=6]{\TFPTalphaBarFiveInvMZPredNum}$ (PDG~2024 $z=\num[round-mode=places,round-precision=3]{\TFPTalphaBarFiveInvMZZ}$);
the inverse bridge back to $\alpha(0)$ is diagnostic and tracked by the explicit error ledger.
The same kernel fixes the axion--photon coupling $g_{a\gamma\gamma}=-4c_3$ and predicts cosmic birefringence,
given a post-recombination excursion $\Delta a=\varphi_0$,
$\beta=\phiz/(4\pi)=\num[round-mode=places,round-precision=4]{\TFPTbetaDegNum}^\circ$, consistent with Planck PR4 observations at the 0.5--1.7$\sigma$ level.
\end{abstract}

%=============================================================================
% KERNEL AND DISCRETE GRAMMAR (Compiler view)
%=============================================================================
\section*{Kernel and Discrete Grammar}
\addcontentsline{toc}{section}{Kernel and Discrete Grammar}

\begin{tcolorbox}[colback=gray!4, colframe=gray!60, title={\textbf{Kernel specification (inputs $\to$ outputs)}}, fonttitle=\bfseries\small]
\textbf{Primitive constant:} $\pi$.

\textbf{Structural input (fixed, not fitted):} the Standard Model abelian trace $\bone=41/10$ (determined by SM particle content and GUT normalization; not a TFPT knob).

\textbf{Notation:} we denote the geometric scale by $\varphi_0$ throughout (suite key: \texttt{varphi0}); no separate $\phi_0$ object is introduced.

\textbf{Kernel invariants (computed):}
\[
c_3=\frac{1}{8\pi},\qquad 
\varphi_{\mathrm{tree}}=\frac{1}{6\pi},\qquad
\delta_{\mathrm{top}}=48c_3^4=\frac{3}{256\pi^4},\qquad
\varphi_0=\varphi_{\mathrm{tree}}+\delta_{\mathrm{top}}.
\]

\textbf{Discrete grammar (no continuous fits):}
\begin{itemize}[leftmargin=1.5em, topsep=2pt, itemsep=1pt]
    \item \textbf{Double cover class:} $k=2$ (exactly two sheets) fixes the backreaction exponent in $\varphi_0(\alpha)$.
    \item \textbf{Holonomy degeneracy:} SU(5) hypercharge holonomy splits into eigenspaces of dimensions $\{3,2\}$, hence $g:=3+2=5$ (a discrete degeneracy label, \emph{not} a gauge coupling).
\end{itemize}

\textbf{Kernel outputs (examples):}
\begin{itemize}[leftmargin=1.5em, topsep=2pt, itemsep=1pt]
    \item $g_{a\gamma\gamma}=-4c_3=-1/(2\pi)$,
    \item $\beta_{\mathrm{rad}}=\varphi_0/(4\pi)$ and $\beta_{\mathrm{deg}}=(180/\pi)\beta_{\mathrm{rad}}$,
    \item two-defect weight $\delta_2=\tfrac{g}{4}(\delta_{\mathrm{top}})^2$ (so $\delta_2/(\delta_{\mathrm{top}})^2=1.25$),
    \item E$_8$ damping anchor $\gamma(0)=\tfrac{g}{g+1}=\tfrac56$.
\end{itemize}

\textbf{Evidence principle:} Every headline number in this paper is either proved from A1--A8 or reproduced by a named verification module in the TFPT suite (artifact locations summarized in Appendix~\ref{app:repro}).
\end{tcolorbox}

\begin{figure}[H]
\centering
\begin{tikzpicture}[
    box/.style={draw, rounded corners, text width=2.55cm, minimum height=1.0cm, align=center, font=\scriptsize},
    arrow/.style={-{Stealth[length=3mm]}, thick}
]
    \node[box, fill=topologyblue!15] (kernel) {Kernel\\$\pi\to(c_3,\varphi_0)$\\SM content: $b_1$\\grammar: $k=2$, $g=5$};
    \node[box, right=0.6cm of kernel, fill=quantumorange!15] (couplings) {Couplings\\CFE $\Rightarrow \alpha$\\policy: $\overline{\alpha}^{(5)}(M_Z)$};
    \node[box, right=0.6cm of couplings, fill=geometrygreen!15] (scales) {Scales\\$R^2$, $f_a$\\$\rho_\Lambda$ candidates};
    \node[box, right=0.6cm of scales, fill=observablepurple!15] (flavor) {Flavor anchors\\$\lambda$, CKM/PMNS};
    \node[box, right=0.6cm of flavor, fill=observablepurple!15] (cosmo) {Cosmology anchors\\$\Omega_b$, $n_s,r$\\$A_s$, $k\to\ell$};

    \draw[arrow] (kernel) -- (couplings);
    \draw[arrow] (couplings) -- (scales);
    \draw[arrow] (scales) -- (flavor);
    \draw[arrow] (flavor) -- (cosmo);
\end{tikzpicture}
\caption{Compiler view: TFPT separates a minimal kernel (discrete invariants) from policy layers (scheme/matching, thresholds, $k\to\ell$ mapping) and produces finite candidate sets rather than continuous fits.}
\label{fig:compiler}
\end{figure}

\begin{table}[H]
\centering
\begin{tabular}{@{}lll@{}}
\toprule
Output & Value & Evidence (module) \\
\midrule
$\alpha^{-1}(0)$ (two-defect, $g=5$) & \num[round-mode=places,round-precision=12]{\TFPTalphaInvGOverFourNum} & \texttt{defect\_partition\_g5\_audit} \\
$\overline{\alpha}^{(5)}(M_Z)^{-1}$ ($\overline{\mathrm{MS}}$ policy) & \num[round-mode=places,round-precision=6]{\TFPTalphaBarFiveInvMZPredNum} & \texttt{alpha\_precision\_audit} \\
$\beta_{\mathrm{deg}}$ (birefringence) & \num[round-mode=places,round-precision=4]{\TFPTbetaDegNum}$^\circ$ & \texttt{birefringence\_tomography} \\
$\Omega_b$ (APS seam; closed v2.8) & \num[round-mode=places,round-precision=6]{\TFPTOmegaBPred} & \texttt{omega\_b\_seam\_embedding\_solver} \\
$\sin^2\theta_{13}$ (TM1/Z$_3$) & \num[round-mode=places,round-precision=6]{\TFPTsinSqThetaThirteen} & \texttt{pmns\_full\_closure\_gate\_v28} \\
$\delta_{\mathrm{CKM}}$ (topology gate; closed v2.8) & $\pi/3$ & \texttt{ckm\_phase\_topology\_referee\_exit\_gate\_v28} \\
$\sin^2\theta_{23}$ (closed v2.8) & \num[round-mode=places,round-precision=6]{\TFPTsinSqThetaTwentyThree} & \texttt{pmns\_full\_closure\_gate\_v28} \\
\bottomrule
\end{tabular}
\caption{Derived constants and anchors from the kernel (v2.8). Each entry is either proved in-text or reproduced by a named suite module; artifact locations are summarized in Appendix~\ref{app:repro}.}
\label{tab:constant-factory}
\end{table}

\begin{figure}[H]
\centering
\begin{tikzpicture}[
    node/.style={draw, rounded corners, align=center, font=\small, minimum width=3.1cm, minimum height=0.9cm},
    arrow/.style={-{Stealth[length=3mm]}, thick}
]
    \node[node, fill=geometrygreen!15] (hol) {SU(5) holonomy\\degeneracy $\{3,2\}$};
    \node[node, right=1.2cm of hol, fill=topologyblue!15] (g) {$g:=3+2=5$\\(discrete node)};
    \node[node, below left=0.9cm and 0.6cm of g, fill=quantumorange!15] (d2) {$\delta_2=\frac{g}{4}(\delta_{\mathrm{top}})^2$\\(two-defect weight)};
    \node[node, below=0.9cm of g, fill=quantumorange!15] (gam) {$\gamma(0)=\frac{g}{g+1}$\\(E$_8$ damping)};
    \node[node, below right=0.9cm and 0.6cm of g, fill=quantumorange!15] (uni) {$\Delta b_3$ candidates include $g/2$\\(unification patch)};

    \draw[arrow] (hol) -- (g);
    \draw[arrow] (g) -- (d2);
    \draw[arrow] (g) -- (gam);
    \draw[arrow] (g) -- (uni);
\end{tikzpicture}
\caption{Crosslink map: the same discrete holonomy degeneracy $g=5$ anchors independent sectors (defect expansion, E$_8$ damping, and a discrete unification patch candidate).}
\label{fig:crosslink}
\end{figure}

\begin{tcolorbox}[colback=blue!3, colframe=blue!40,
    title={\textbf{Origin of $g=5$: Holonomy Spectrum of the Hypercharge Operator}},
    fonttitle=\bfseries\small]
The holonomy degeneracy $g=5$ is not a fit parameter but the dimension of the
SU(5) fundamental representation in which the Standard Model hypercharge
generator $Y=\mathrm{diag}(-\tfrac13,-\tfrac13,-\tfrac13,\tfrac12,\tfrac12)$ acts.
The eigenvalue spectrum decomposes as $\{3,2\}$ (multiplicities of
$Y=-1/3$ and $Y=+1/2$), giving $g=3+2=5$.

\textbf{Rigidity:} perturbations of $Y$ away from the exact SM charges
break the integer split immediately.
\textbf{Evidence:} \texttt{holonomy\_split\_from\_hypercharge\_trace\_audit}.
\end{tcolorbox}

%=============================================================================
% CLAIM MAP - Reader's Contract
%=============================================================================
\vspace{1em}
\begin{tcolorbox}[colback=gray!5, colframe=gray!60, title={\textbf{Claim Map: What This Paper Establishes}}, fonttitle=\bfseries]
\textbf{Theorems (given axioms A1--A8):}
\begin{itemize}[leftmargin=1.5em, topsep=2pt, itemsep=1pt]
    \item Fine-structure constant (CFE closure): one-defect truncation gives $\alpha^{-1}(0)=\num[round-mode=places,round-precision=10]{\TFPTalphaInvSelfNum}$ ($-0.037$ ppm vs CODATA~2022). A two-defect refinement anchored by holonomy degeneracy $g=5$ gives $\alpha^{-1}(0)=\num[round-mode=places,round-precision=12]{\TFPTalphaInvGOverFourNum}$ (CODATA $z=\num[round-mode=places,round-precision=3]{\TFPTalphaZGOverFourNum}$).
    \item Cosmic birefringence: $\beta = \num[round-mode=places,round-precision=4]{\TFPTbetaDegNum}^\circ$ (UFE + minimal topological excursion, $n=1$)
    \item Cabibbo angle: $\lambda = \num[round-mode=places,round-precision=6]{0.22437623688472177}$ (derived from $\phiz$; sub-permille deviation vs reference)
    \item Two-loop RG fingerprints: $\alpha_3(1\,\mathrm{PeV})=\num[round-mode=places,round-precision=6]{\TFPTalphaThreePeVNum}$ (rel.\ dev.\ vs.\ $\varphi_0$: $\num[round-mode=places,round-precision=3]{\TFPTalphaThreePeVRelDevPercent}\%$) and exact crossing scales $\alpha_3(\mu)=\varphi_0$, $\alpha_3(\mu)=c_3$ (checked by \texttt{two\_loop\_rg\_fingerprints} in the verification suite)
\end{itemize}

\textbf{Structural extensions (status: testable):}
\begin{itemize}[leftmargin=1.5em, topsep=2pt, itemsep=1pt]
    \item E$_8$ cascade with $\gamma(0)=g/(g+1)=5/6$ (anchored by holonomy degeneracy $g=5$)
    \item Block constants $\zeta_B$ with discrete indices $k_B, r_B$
    \item Z$_3$ flavor architecture with single phase $\delta$
    \item PMNS: $\sin^2\theta_{13} \approx 0.0231$ via TM1 pattern (few-\% dev.)
    \item Inflation: $A_s \sim 2\times10^{-9}$, $n_s \approx 0.964$, $r \approx 0.0038$ via $R^2$ completion (derived under assumptions K1--K3; reproduced by the verification module \texttt{effective\_action\_r2} in the suite)
    \item Bounce transfer function $T(k)$: large-scale features from torsion-driven bounce (conditional; evaluated by the verification module \texttt{bounce\_perturbations} in the suite)
    \item Axion DM: $m_a \approx 65.19\,\mu$eV, $\nu \approx 15.76$ GHz (background + DM modes)
\end{itemize}

\textbf{Direct falsification paths:}
\begin{itemize}[leftmargin=1.5em, topsep=2pt, itemsep=1pt]
    \item Any future $\alpha$ measurement deviating $>1$ ppm from TFPT
    \item Birefringence $\beta$ outside $[0.1^\circ, 0.4^\circ]$ with $<0.05^\circ$ error
    \item RG fingerprint $\alpha_3(1\,\mathrm{PeV})$ deviating $>2\%$ from $\phiz$ (see \texttt{two\_loop\_rg\_fingerprints})
    \item Inflation $r$ outside $[0.001, 0.005]$ with precision $\sigma_r < 0.001$
    \item PMNS $\sin^2\theta_{13}$ deviating $>10\%$ from $0.0231$
    \item Axion haloscope non-detection at $15.6 \pm 0.5$ GHz (given local DM density)
\end{itemize}
\end{tcolorbox}

%=============================================================================
% ASSUMPTION LEDGER - Transparency about what is assumed
%=============================================================================
\vspace{0.5em}
\begin{tcolorbox}[colback=yellow!5, colframe=yellow!60!black, 
    title={\textbf{Assumption Ledger: Established vs.\ New}}, fonttitle=\bfseries]

\textbf{Established physics (citable):}
\begin{itemize}[leftmargin=1.5em, topsep=2pt, itemsep=1pt]
    \item Riemann--Cartan geometry with torsion (Hehl et al.~1976)
    \item Fujikawa anomaly and chiral measure (Fujikawa 1979)
    \item Background-field gauge and Ward identities (standard QFT)
    \item Maxwell convention: U(1) field strength defined as $F=dA$ (gauge invariance exact in a Riemann--Cartan background)
    \item SM abelian trace $\bone = 41/10$ (particle content, GUT normalization)
    \item GR limit requirement ($K \to 0 \Rightarrow$ Einstein--Hilbert)
\end{itemize}

\textbf{Structural inputs (to be justified or tested):}
\begin{center}
\begin{tabular}{@{}p{4cm}p{4cm}p{4.5cm}@{}}
\toprule
\textbf{Input} & \textbf{Justification} & \textbf{Falsification Test}\\
\midrule
Orientable double cover (A2) & Spinor structure requires orientability & CPT tests, spin-torsion bounds\\
$\phitree = 1/(6\pi)$ from $\mathbb{Z}_2$ gluing term (A3) & $\eta$-gluing on cut double cover & Change gluing class/topology $\Rightarrow$ shift $\alpha$\\
$\deltatop = 48\cthree^4$ & Spin-lifted $\mathbb{Z}_2$ deficit factor & See Appendix~\ref{app:48deriv}\\
Backreaction exponent $=2$ & Defect-sector reweighting on the orientable double cover (Lemma~\ref{lem:defectreweight}) & Sensitivity analysis (Sec.~\ref{sec:sensitivity})\\
Holonomy degeneracy $g=5$ & SU(5) hypercharge holonomy eigenspaces $\{3,2\}$ (kernel crosslink) & Crosslink failures in $\{\delta_2,\gamma(0),$ unification patch$\}$\\
$\delta_2=\frac54\,(\deltatop)^2$ (two-defect) & Defect partition enumeration anchored by $g=5$ (theoryv3 audit) & $\alpha^{-1}(0)$ shifts outside CODATA window\\
$\Delta a_{\text{top}} = n\,\phiz$ (birefringence) & Topological excursion quantization (Definition~\ref{def:topoexcursion}); Sec.~\ref{sec:deltaa} gives one minimal dynamical realization & $\beta(z)$ tomography inconsistent\\
$\gamma(0)=g/(g+1)=5/6$ (E$_8$) & Anchored by holonomy degeneracy $g=5$ (crosslink) & Block-scale deviations $>5\%$\\
$R^2$ completion with $M/\Mpl=\sqrt{8\pi}\cthree^4$ & Derived under assumptions K1--K3 (Appendix~\ref{app:r2}); reproduced by \texttt{effective\_action\_r2} in the verification suite & $(A_s,n_s,r)$ inconsistent\\
\bottomrule
\end{tabular}
\end{center}

\textbf{Parameter count:} The core CFE uses 0 free parameters. With $\gamma(0)=g/(g+1)=5/6$, the E$_8$ cascade 
introduces no continuous fit parameters (only discrete E$_8$ data and the TFPT invariants). 
The inputs above are \emph{structural choices}, not continuous fits.
\end{tcolorbox}

\tableofcontents
%=============================================================================
\section{Introduction: The Problem of Parameters}
%=============================================================================

The Standard Model of particle physics contains approximately 19 free parameters 
that must be determined experimentally. Among these, the fine-structure constant 
$\alpha \approx 1/137$ stands out as a dimensionless number whose value has no 
explanation within conventional frameworks.

TFPT addresses this problem by demonstrating that $\alpha$ is not a free parameter 
but a \emph{fixed point} of an effective potential, uniquely determined by:
\begin{enumerate}[label=(\roman*)]
    \item Topological quantization from 11D Chern--Simons theory
    \item Geometric reduction on a M\"obius fiber
    \item Quantum consistency via background-field gauge Ward identities
\end{enumerate}

\begin{figure}[H]
\centering
\begin{tikzpicture}[
    scale=1.0,
    box/.style={draw, rounded corners, minimum width=2.8cm, minimum height=1cm, align=center, font=\small},
    arrow/.style={-{Stealth[length=3mm]}, thick}
]
    % Boxes
    \node[box, fill=topologyblue!20] (topo) at (0,3) {Topology\\$\cthree = \frac{1}{8\pi}$};
    \node[box, fill=geometrygreen!20] (geom) at (4,3) {Geometry\\$\phiz = \frac{1}{6\pi}+\frac{3}{256\pi^4}$};
    \node[box, fill=quantumorange!20] (qft) at (2,1.5) {Quantum FPE\\CFE: $\alpha^3 - 2\cthree^3\alpha^2 - \ldots = 0$};
    \node[box, fill=observablepurple!20] (obs1) at (-0.5,0) {$\alpha^{-1}(0)$\\$\num[round-mode=places,round-precision=4]{\TFPTalphaInvGOverFourNum}$};
    \node[box, fill=observablepurple!20] (obs2) at (2,0) {$\gagg$\\$-4\cthree$};
    \node[box, fill=observablepurple!20] (obs3) at (4.5,0) {$\beta$\\$\num[round-mode=places,round-precision=4]{\TFPTbetaDegNum}^\circ$};
    
    % Arrows
    \draw[arrow, topologyblue] (topo) -- (qft);
    \draw[arrow, geometrygreen] (geom) -- (qft);
    \draw[arrow, observablepurple] (qft) -- (obs1);
    \draw[arrow, observablepurple] (qft) -- (obs2);
    \draw[arrow, observablepurple] (qft) -- (obs3);
    
    % Labels
    \node[font=\footnotesize, topologyblue] at (-0.8,2.3) {fixes $A$};
    \node[font=\footnotesize, geometrygreen] at (4.8,2.3) {fixes $\mathcal{K}$};
\end{tikzpicture}
\caption{Logical flow of TFPT: topological and geometric invariants determine quantum 
fixed-point parameters, which yield observables without free inputs.}
\label{fig:flow}
\end{figure}

%=============================================================================
\section{Axiomatic Foundation}
%=============================================================================

TFPT is built on eight axioms, organized into three categories: geometric structure, 
quantum field theory, and closure conditions. Each axiom is either derived from 
established physics or explicitly marked as a new postulate.

%-----------------------------------------------------------------------------
